\begin{musica}{Assim Mesmo É Que É (Rapariga)}{António Vicente}{Estudantina Universitária de Coimbra}
\nota{Instrumental: F C G C (x2)}

\ac{G}{}Lá da aldeia d'onde sou
\ac{C}{}Não perdoo as raparigas
\ac{G}{}Se uma o olho me piscou
\ac{C}{}Mete-mo logo em intrigas

\ac{F}{}Dou-lhe dois ou três beijinhos
\ac{C}{}E vai de bater o pé
\ac{G}{}Eu não quero mexericos
\ac{C}{}E assim mesmo é que é
\ac{F G}{}Que eu não quero mexericos
\ac{C}{}E assim mesmo é que é

\nota{Refrão:}
\ac{F C}{}Ai rapariga, se fores à fonte
\ac{G C}{}Vai pelo carreiro que chegas lá mais depressa
\ac{F C}{}Ai tem cuidado com os rapazes
\ac{G C}{}Loucos por ti, vê lá se algum tropeça

\nota{Instrumental}

No outro dia a Rosita
Que é baixinha e trigueira
Foi ao baile com o António
Andaram em brincadeiras

E agora já namoram
E é tão bom de ver ai é
Qualquer dia hão de casar
E assim mesmo é que é
Qualquer dia hão de casar
E assim mesmo é que é

\nota{Refrão}

\nota{Instrumental}

Esta vida são dois dias
Diz o povo e tem razão
Se é tão competente
Vou gozá-la até mais não

E se encontrar a minha amada
Sorridente cheio de fé
Vou levá-la ao altar
E assim mesmo é que é
Vou levá-la ao altar
E assim mesmo é que é

\nota{Refrão}

\nota{Instrumental}

\ac{F C}{}Ai rapariga, rapariga, rapariga
\ac{G C}{}Rapariga, ai rapariga tem cuidado
\ac{F C}{}Ai rapariga, rapariga, rapariga
\ac{G C}{}Rapariga, rapariga e assim mesmo é que é

\nota{G C}

\posicoes{C: x32010, F: 133211, G: 320033}
\end{musica}
